% ============================================================
% CAPITOLO 34 - LA MODULAZIONE AI TONI VICINI: TEORIA
% ============================================================
\chapter{La modulazione ai toni vicini: teoria}
\label{chap:modulazione}

Finora si è lavorato dentro un'unica tonalità. La \textbf{modulazione} è il passaggio da una tonalità a un'altra. È un processo graduale che sfrutta le relazioni fra le tonalità e la capacità di certi accordi di avere un doppio significato. Questo capitolo ne dà la teoria per i \textbf{toni vicini}, cioè le tonalità che condividono con quella di partenza la maggior parte delle note; la tecnica è nel Capitolo \ref{chap:tecnica_modulazione}.

\section{Che cos'è una modulazione}
\label{sec:definizione_modulazione}

\begin{definizione}[Modulazione]
	\label{def:modulazione}
	Una \textbf{modulazione} è il processo con cui una composizione abbandona stabilmente la tonalità di partenza (tonalità d'impianto) per stabilirne una nuova. Comporta:
	\begin{enumerate}
		\item l'affermazione chiara della tonalità di partenza;
		\item un percorso armonico di transizione;
		\item la conferma della nuova tonalità con una cadenza.
	\end{enumerate}
	Si distingue dall'\textbf{accenno} o dalla \textbf{tonicizzazione} (Capitolo \ref{chap:tonicizzazione}), in cui si crea per un momento l'illusione di una nuova tonalità senza lasciare stabilmente quella di partenza.
\end{definizione}

Nella tonicizzazione la nuova tonalità viene «visitata» brevemente, senza una cadenza autentica a confermarla. Nella modulazione il cambiamento è strutturale e la nuova tonalità diventa il nuovo centro di gravità.

\section{Le tre fasi del processo modulante}
\label{sec:fasi_modulazione}

\begin{regola}[Le tre fasi della modulazione]
	\label{reg:tre_fasi_modulazione}
	Il processo modulante si sviluppa in tre momenti:
	\begin{enumerate}
		\item \textbf{Affermazione (o consolidamento):} la tonalità di partenza è stabilita con chiarezza da progressioni caratteristiche e, preferibilmente, da una cadenza. L'ascoltatore deve percepire senza ambiguità il centro tonale da cui si parte.

		\item \textbf{Transizione:} si introduce un \textbf{accordo pivot} (o cardine), che appartiene a entrambe le tonalità. Funge da ponte armonico fra la tonalità vecchia e la nuova, senza discontinuità.

		\item \textbf{Conferma (o cadenza):} la nuova tonalità è confermata da una progressione cadenzale chiara, di norma una cadenza perfetta (V--I), spesso preceduta dalla sottodominante (II--V--I, IV--V--I). La fase è essenziale: senza una cadenza convincente nella nuova tonalità si ha solo una tonicizzazione passeggera.
	\end{enumerate}
\end{regola}

Lo schema in tre tempi è un modello. Nelle composizioni brevi l'affermazione può essere implicita e la transizione molto rapida. Non può mai mancare la \textbf{conferma} della nuova tonalità.

\section{I toni vicini}
\label{sec:toni_vicini}

Le tonalità non sono tutte «lontane» allo stesso modo: alcune condividono molte note, altre poche. Questa distanza si legge sul circolo delle quinte (Figura \ref{fig:circolo_quinte_migliorato_compatto}).

\begin{definizione}[Tonalità vicine]
	\label{def:tonalita_vicine}
	Due tonalità sono \textbf{vicine} se le armature di chiave differiscono al massimo di \textbf{un'alterazione} (un diesis o un bemolle). La vicinanza significa che condividono sei note su sette della scala. Sul circolo delle quinte sono vicine le tonalità adiacenti o legate dal rapporto di relativa maggiore/minore.
\end{definizione}

\subsection{I cinque toni vicini di una tonalità maggiore}
\label{subsec:toni_vicini_maggiore}

\begin{regola}[Tonalità vicine al modo maggiore]
	\label{reg:toni_vicini_maggiore}
	Data una tonalità maggiore, i suoi cinque toni vicini sono:
	\begin{enumerate}
		\item la \textbf{relativa minore} (stessa armatura);
		\item la tonalità della \textbf{dominante} (una quinta sopra: un diesis in più o un bemolle in meno);
		\item la \textbf{relativa minore della dominante};
		\item la tonalità della \textbf{sottodominante} (una quinta sotto: un bemolle in più o un diesis in meno);
		\item la \textbf{relativa minore della sottodominante}.
	\end{enumerate}
\end{regola}

\begin{esempio}[I cinque toni vicini di Do maggiore]
	\label{ex:toni_vicini_do}
	Da Do maggiore (nessuna alterazione) i toni vicini sono:
	\begin{itemize}
		\item La minore (relativa, nessuna alterazione);
		\item Sol maggiore (dominante, un diesis);
		\item Mi minore (relativa di Sol, un diesis);
		\item Fa maggiore (sottodominante, un bemolle);
		\item Re minore (relativa di Fa, un bemolle).
	\end{itemize}
	Tutte condividono sei note su sette con Do maggiore.
\end{esempio}

\subsection{I cinque toni vicini di una tonalità minore}
\label{subsec:toni_vicini_minore}

\begin{regola}[Tonalità vicine al modo minore]
	\label{reg:toni_vicini_minore}
	Data una tonalità minore, i suoi cinque toni vicini sono:
	\begin{enumerate}
		\item la \textbf{relativa maggiore} (stessa armatura della minore naturale), cioè la tonalità del III grado;
		\item la tonalità del \textbf{IV grado} (sottodominante);
		\item la tonalità del \textbf{V grado} (dominante);
		\item la tonalità del \textbf{VI grado} (sopradominante);
		\item la tonalità del \textbf{VII grado}.
	\end{enumerate}
	Nel minore la situazione è un po' più complessa per le varianti armonica e melodica della scala, ma il criterio della vicinanza per alterazioni resta valido.
\end{regola}

\begin{esempio}[I cinque toni vicini di La minore]
	\label{ex:toni_vicini_la_m}
	Da La minore (nessuna alterazione in chiave) i toni vicini sono:
	\begin{itemize}
		\item Do maggiore (relativa, nessuna alterazione);
		\item Mi minore (dominante, un diesis);
		\item Re minore (sottodominante, un bemolle);
		\item Sol maggiore (VII grado, dominante della relativa, un diesis);
		\item Fa maggiore (VI grado, sottodominante della relativa, un bemolle).
	\end{itemize}
\end{esempio}

\section{L'accordo pivot}
\label{sec:accordo_pivot}

La modulazione ai toni vicini si basa sugli \textbf{accordi pivot} (cardine, di perno).

\begin{definizione}[Accordo pivot]
	\label{def:accordo_pivot}
	Un \textbf{accordo pivot} appartiene contemporaneamente a due tonalità e ha in ciascuna una funzione armonica precisa. Nella modulazione viene prima percepito come parte della tonalità di partenza, poi «reinterpretato» come parte di quella di arrivo.
\end{definizione}

Non tutti gli accordi sono ugualmente efficaci come pivot. Contano due fattori:

\begin{enumerate}
	\item \textbf{la condivisione delle note:} l'accordo deve essere costruibile con le note di entrambe le scale;
	\item \textbf{la forza funzionale nella nuova tonalità:} un pivot che nella nuova tonalità è II, IV o V grado è più efficace di uno che è III o VI.
\end{enumerate}

\subsection{Accordi pivot tra tonalità relative}
\label{subsec:pivot_relative}

La modulazione tra tonalità relative (Do maggiore $\leftrightarrow$ La minore) è la più semplice, perché condividono tutte le note della scala.

\begin{regola}[Accordi pivot tra relative]
	\label{reg:pivot_relative}
	Tra Do maggiore e La minore gli accordi pivot più efficaci sono:
	\begin{itemize}
		\item \textbf{Re minore:} II grado in Do M, IV grado in La m (sottodominante in entrambe);
		\item \textbf{Mi minore:} III grado in Do M, V grado in La m (da tonica a dominante);
		\item \textbf{Fa maggiore:} IV grado in Do M, VI grado in La m.
	\end{itemize}
	Il più efficace è di solito il II di Do M / IV di La m, perché ha funzione di sottodominante forte in entrambe.
\end{regola}

\begin{esempio}[Modulazione Do maggiore $\rightarrow$ La minore con pivot Re minore]
	\label{ex:mod_do_lam}
	Il Re minore (II di Do M) è reinterpretato come IV di La minore e prepara la cadenza V--I in La minore.
	\musica{es_cap16_01}
\end{esempio}

\subsection{Accordi pivot verso la dominante}
\label{subsec:pivot_dominante}

La modulazione dalla tonica alla sua dominante (Do M $\rightarrow$ Sol M) è altrettanto naturale.

\begin{regola}[Accordi pivot verso la dominante maggiore]
	\label{reg:pivot_dominante}
	Tra Do maggiore e Sol maggiore gli accordi pivot più comuni sono:
	\begin{itemize}
		\item \textbf{Do maggiore stesso:} I grado in Do M, IV grado in Sol M;
		\item \textbf{Mi minore:} III grado in Do M, VI grado in Sol M;
		\item \textbf{Sol maggiore:} V in Do M, I in Sol M (meno efficace come pivot, perché in Do suona già come dominante; la \sigla{V}{7}{} di Do, con il Fa naturale, non appartiene a Sol M).
	\end{itemize}
	Il più efficace è il Do maggiore reinterpretato come IV di Sol, perché permette di far seguire subito la cadenza V--I nella nuova tonalità.
\end{regola}

\begin{esempio}[Modulazione Do maggiore $\rightarrow$ Sol maggiore con pivot I = IV]
	\label{ex:mod_do_sol}
	L'accordo di Do maggiore, sentito all'inizio come tonica, diventa sottodominante di Sol maggiore.
	\musica{es_cap16_02}
\end{esempio}

\subsection{Accordi pivot verso la sottodominante}
\label{subsec:pivot_sottodominante}

\begin{regola}[Accordi pivot verso la sottodominante maggiore]
	\label{reg:pivot_sottodominante}
	Tra Do maggiore e Fa maggiore:
	\begin{itemize}
		\item \textbf{Fa maggiore stesso:} IV in Do M, I in Fa M;
		\item \textbf{Re minore:} II in Do M, VI in Fa M;
		\item \textbf{Do maggiore:} I in Do M, V in Fa M.
	\end{itemize}
	Il più efficace è il Do maggiore reinterpretato come dominante di Fa, che porta direttamente alla cadenza.
\end{regola}

\begin{esempio}[Modulazione Do maggiore $\rightarrow$ Fa maggiore con pivot I = V]
	\label{ex:mod_do_fa}
	\musica{es_cap16_03}
\end{esempio}

\section{Dalla tonicizzazione alla modulazione}
\label{sec:collegamento_mod}

Tonicizzazione e modulazione sono due punti di un \textit{continuum}. Una tonicizzazione prolungata e confermata da una cadenza diventa a tutti gli effetti una modulazione.
